\newpage
\section{Experimentos y Análisis de Resultados}
\label{sec:experimentos}

\subsection{Casos, Parámetros y Metodología}

Configuración experimental utilizada:

\begin{itemize}[noitemsep]
    \item 3 mercados: IBEX 35, S\&P 100 y S\&P 500.
    \item 50 ejecuciones independientes por algoritmo estocástico.
    \item Presupuesto de 10.000 evaluaciones por ejecución.
    \item Semilla base 42, incrementándose secuencialmente ($42+i$).
    \item Entrenamiento en 2015--2024 y validación en 2025.
\end{itemize}

Todos estos parámetros se leen desde \texttt{config.cfg}, por lo que no es
necesario recompilar para modificarlos.

\subsection{Salida Resumida de los Algoritmos Obligatorios}

Las tablas siguientes recogen la salida media de las ejecuciones para los
algoritmos obligatorios de la práctica. Se muestran solo estos métodos para que
la comparación quede centrada en la parte común del enunciado.

\begin{table}[H]
\centering
\scriptsize
\begin{tabular}{@{}lrrrrrr@{}}
\toprule
\textbf{IBEX 35} & \textbf{2015-2024 Fitness} & \textbf{2025 Fitness} & \textbf{Profit} & \textbf{Desv} & \textbf{Evals} & \textbf{Tiempo(s)} \\
\midrule
Greedy    & -4.794 & -3.227 & -0.033 & 0.000 & 1     & 0.000 \\
Random    & -4.911 & -3.305 &  0.052 & 0.091 & 10000 & 0.024 \\
BL        & -4.614 & -3.116 &  0.021 & 0.039 & 1078  & 0.001 \\
AGG-Arit  & -4.663 & -3.142 &  0.042 & 0.052 & 10000 & 0.009 \\
AGG-BLX   & -4.496 & -3.091 &  0.036 & 0.006 & 10000 & 0.010 \\
AGE-Arit  & -4.587 & -3.119 &  0.043 & 0.039 & 10000 & 0.010 \\
AGE-BLX   & -4.499 & -3.097 &  0.033 & 0.010 & 10000 & 0.011 \\
AM-All    & -4.591 & -3.118 &  0.036 & 0.035 & 10000 & 0.008 \\
AM-Rand   & -4.490 & -3.084 &  0.033 & 0.000 & 10000 & 0.009 \\
AM-Best   & -4.490 & -3.083 &  0.033 & 0.000 & 10000 & 0.009 \\
\bottomrule
\end{tabular}
\caption{Resultados medios en IBEX 35 para los algoritmos obligatorios.}
\label{tab:ibex-oblig}
\end{table}

\begin{table}[H]
\centering
\scriptsize
\begin{tabular}{@{}lrrrrrr@{}}
\toprule
\textbf{S\&P 100} & \textbf{2015-2024 Fitness} & \textbf{2025 Fitness} & \textbf{Profit} & \textbf{Desv} & \textbf{Evals} & \textbf{Tiempo(s)} \\
\midrule
Greedy    & -3.305 & -3.791 & 0.100 & 0.000 & 1     & 0.000 \\
Random    & -3.400 & -4.080 & 0.079 & 0.041 & 10000 & 0.098 \\
BL        & -3.116 & -3.890 & 0.097 & 0.084 & 4351  & 0.033 \\
AGG-Arit  & -3.120 & -3.845 & 0.083 & 0.041 & 10000 & 0.071 \\
AGG-BLX   & -2.879 & -3.721 & 0.096 & 0.023 & 10000 & 0.075 \\
AGE-Arit  & -3.157 & -3.872 & 0.082 & 0.054 & 10000 & 0.071 \\
AGE-BLX   & -2.937 & -3.744 & 0.094 & 0.037 & 10000 & 0.075 \\
AM-All    & -3.110 & -3.830 & 0.085 & 0.037 & 10000 & 0.073 \\
AM-Rand   & -2.822 & -3.735 & 0.099 & 0.011 & 10000 & 0.074 \\
AM-Best   & -2.814 & -3.740 & 0.102 & 0.008 & 10000 & 0.074 \\
\bottomrule
\end{tabular}
\caption{Resultados medios en S\&P 100 para los algoritmos obligatorios.}
\label{tab:sp100-oblig}
\end{table}

\begin{table}[H]
\centering
\scriptsize
\begin{tabular}{@{}lrrrrrr@{}}
\toprule
\textbf{S\&P 500} & \textbf{2015-2024 Fitness} & \textbf{2025 Fitness} & \textbf{Profit} & \textbf{Desv} & \textbf{Evals} & \textbf{Tiempo(s)} \\
\midrule
Greedy    & -3.370 & -3.909 & 0.024 & 0.000 & 1     & 0.000 \\
Random    & -3.648 & -4.364 & 0.019 & 0.044 & 10000 & 1.851 \\
BL        & -3.428 & -4.330 & 0.004 & 0.154 & 9645  & 1.526 \\
AGG-Arit  & -2.961 & -3.958 & 0.018 & 0.082 & 10000 & 1.504 \\
AGG-BLX   & -2.791 & -3.893 & 0.023 & 0.046 & 10000 & 1.490 \\
AGE-Arit  & -3.176 & -4.073 & 0.018 & 0.074 & 10000 & 1.478 \\
AGE-BLX   & -3.012 & -3.984 & 0.023 & 0.066 & 10000 & 1.483 \\
AM-All    & -3.204 & -4.102 & 0.015 & 0.040 & 10000 & 1.573 \\
AM-Rand   & -2.738 & -3.907 & 0.017 & 0.034 & 10000 & 1.537 \\
AM-Best   & -2.731 & -3.905 & 0.017 & 0.035 & 10000 & 1.538 \\
\bottomrule
\end{tabular}
\caption{Resultados medios en S\&P 500 para los algoritmos obligatorios.}
\label{tab:sp500-oblig}
\end{table}

\subsection{Resultados Obtenidos (Algoritmos Obligatorios)}

La salida completa de ejecución muestra que, en términos de \textit{train}, los
métodos evolutivos y meméticos mejoran de forma clara al Greedy, pero la mejora
se reduce cuando se mide en \textit{test}. Esa diferencia entre entrenamiento y
validación es precisamente la que aparece reflejada en los boxplots: algunas
estrategias alcanzan soluciones muy buenas en promedio, pero con mayor
dispersión y presencia de valores atípicos.

Si se toma Greedy como referencia\footnote{$\text{Mejora(\%)} = \frac{|\text{Algoritmo} - \text{Greedy}|}{|\text{Greedy}|} \times 100$, donde los valores provienen de la tabla de resultados.}, la mejora media en \textit{train} de las
mejores configuraciones ronda aproximadamente un \(6\%\) en IBEX 35, un
\(15\%\) en S\&P 100 y un \(19\%\) en S\&P 500. En \textit{test}, en cambio,
la ganancia se comprime: alrededor de un \(4\%-5\%\) en IBEX 35, cerca de un
\(2\%\) en S\&P 100 y apenas un \(1\%\) en S\&P 500. Esto indica que el
problema se vuelve más exigente al crecer el mercado y que el sobreajuste entre
entrenamiento y validación gana peso en los casos grandes.

\subsection{Análisis de las Distribuciones}

Las Figuras~\ref{fig:ibex-box}, \ref{fig:sp100-box} y \ref{fig:sp500-box}
permiten observar la estabilidad de cada algoritmo. Cada mercado se muestra en
dos partes para mantener la legibilidad. Greedy aparece como una raya
horizontal porque no tiene variabilidad: sirve como referencia rápida, pero
queda lejos de los métodos poblacionales en \textit{train}. Random es el más
irregular en todos los mercados, lo que podría confirmanos que la aleatoriedad es impredecible.

En los algoritmos genéticos, el cruce BLX domina de forma bastante clara al
aritmético. La mejora visual no es solo de posición de la mediana: también se
observa una reducción de dispersión y menos ejecuciones extremas. En términos
prácticos, BLX suele aportar una ganancia adicional de entre un \(1\%\) y un
\(3\%\) en \textit{test} frente al cruce aritmético, con una estabilidad mayor.
Esto se aprecia especialmente en IBEX 35 y S\&P 100, donde las cajas de BLX son
más compactas y los whiskers quedan mejor controlados.

La comparación entre AGG y AGE muestra que no existe una superioridad universal
del esquema estacionario sobre el generacional. En IBEX 35, ambos se mueven en
rangos parecidos, aunque AGE no aporta una mejora clara sobre AGG. En S\&P 100,
AGG-BLX y AGE-BLX quedan muy cerca, pero AGG suele mantener una mediana algo
mejor y menos ruido. En S\&P 500, la tendencia se mantiene: el modelo generacional
con BLX ofrece una combinación más equilibrada entre calidad media y robustez.

Los boxplots también muestran el papel de los valores atípicos en las distintas ejecuciones. En casi todos los gráficos aparecen uno o dos outliers aislados, lo que indica que incluso cuando la dispersión central es reducida todavía existen ejecuciones puntuales alejadas del comportamiento habitual. 

Entre los algoritmos obligatorios, AM-All presenta una dispersión superior a AM-Rand y AM-Best. Esa diferencia confirma que aplicar la búsqueda local a toda la población incrementa el coste de explotación sin traducirse en una mejora proporcional.

Entre las variantes meméticas, AM-Rand y AM-Best quedan prácticamente empatadas,
con cajas estrechas y mediana alta. En varios mercados apenas se distinguen en
los gráficos, lo que indica que la selección aleatoria de una fracción de la
población o la selección del bloque más prometedor terminan convergiendo a una
calidad muy similar. En cambio, AM-All suele quedar por detrás y es más sensible
a ejecuciones malas, por lo que no compensa su mayor intensidad de búsqueda
local.

\subsection{Resumen Comparativo}

Si se resume el comportamiento global de los gráficos, pueden extraerse cuatro
ideas principales. Primero, el rendimiento relativo empeora al pasar de IBEX 35
a S\&P 100 y, sobre todo, a S\&P 500: el problema grande es más difícil y las
ventajas sobre Greedy se concentran en \textit{train}, no tanto en \textit{test}.
Segundo, BLX es la opción más robusta frente al cruce aritmético, tanto en la
calidad central como en la dispersión. Tercero, la búsqueda local sobre toda la
población no resulta la mejor estrategia; es preferible limitarla a una parte de
los individuos o concentrarla en los mejores candidatos. Cuarto, los algoritmos
con mejores medianas no siempre son los que presentan menos outliers, por lo que
la estabilidad debe considerarse junto con el valor medio.

En cuanto al coste temporal, la diferencia entre algoritmos es pequeña frente al
presupuesto total de evaluaciones. Los métodos más simples apenas tardan unas
pocas centésimas de segundo en IBEX y S\&P 100, mientras que en S\&P 500 el
tiempo aumenta, pero sigue estando dominado por el coste de evaluación y no por
la lógica concreta del operador evolutivo. Por tanto, la selección del algoritmo
debe hacerse más por calidad y estabilidad que por tiempo de ejecución.

\begin{figure}[H]
    \centering
    \begin{minipage}[t]{0.9\textwidth}
        \centering
        \includegraphics[width=\textwidth]{figuras/fig_ibex_oblig_pt1.png}

        \smallskip
        {\small Parte 1}
    \end{minipage}
    
    \vspace{0.5cm}
    
    \begin{minipage}[t]{0.9\textwidth}
        \centering
        \includegraphics[width=\textwidth]{figuras/fig_ibex_oblig_pt2.png}

        \smallskip
        {\small Parte 2}
    \end{minipage}
    \caption{Distribución del \textit{fitness} por algoritmo en IBEX 35.}
    \label{fig:ibex-box}
\end{figure}

\begin{figure}[H]
    \centering
    \begin{minipage}[t]{0.9\textwidth}
        \centering
        \includegraphics[width=\textwidth]{figuras/fig_sp100_oblig_pt1.png}

        \smallskip
        {\small Parte 1}
    \end{minipage}
    
    \vspace{0.5cm}
    
    \begin{minipage}[t]{0.9\textwidth}
        \centering
        \includegraphics[width=\textwidth]{figuras/fig_sp100_oblig_pt2.png}

        \smallskip
        {\small Parte 2}
    \end{minipage}
    \caption{Distribución del \textit{fitness} por algoritmo en S\&P 100.}
    \label{fig:sp100-box}
\end{figure}

\begin{figure}[H]
    \centering
    \begin{minipage}[t]{0.9\textwidth}
        \centering
        \includegraphics[width=\textwidth]{figuras/fig_sp500_oblig_pt1.png}

        \smallskip
        {\small Parte 1}
    \end{minipage}
    
    \vspace{0.5cm}
    
    \begin{minipage}[t]{0.9\textwidth}
        \centering
        \includegraphics[width=\textwidth]{figuras/fig_sp500_oblig_pt2.png}

        \smallskip
        {\small Parte 2}
    \end{minipage}
    \caption{Distribución del \textit{fitness} por algoritmo en S\&P 500.}
    \label{fig:sp500-box}
\end{figure}
